
\documentclass[11pt]{article}
\usepackage{fullpage}
\usepackage[left=1in,top=1in,right=1in,bottom=1in,headheight=3ex,headsep=3ex]{geometry}
\usepackage{graphicx}
\usepackage{float}
\usepackage{enumerate}
\usepackage{sgamevar, tikz} % Game theory packages
\usepackage{pgfplots}

\usepackage{sgamevar, tikz} % Game theory packages
\usetikzlibrary{trees, calc} % For extensive form games

\newcommand{\blankline}{\quad\pagebreak[2]}
%%%%%%%%%%%%%%%%%%%%%%%%%%%%%%%%%%%%%%%%%%%%%%%%%%%%%%%%%%%%%%

% Modify Course title, instructor name, semester here %%%%%%%%

\title{Microeconom\'ia Aplicada (MAE)}
\author{Ayudantia 2: Juegos Din\'amicos con informaci\'on completa \\ Profesor: Jose Manuel Paz y Mi\~no \\ Ayudante: Vicente Merrill}
\date{}

% Don't touch this %%%%%%%%%%%%%%%%%%%%%%%%%%%%%%%%%%%%%%%%%%%
\usepackage[sc]{mathpazo}
\linespread{1.05} % Palatino needs more leading (space between lines)
\usepackage[T1]{fontenc}
\usepackage{setspace}
\usepackage{multicol}
\usepackage{fancyhdr,lastpage}
\usepackage{url}
\pagestyle{fancy}
\usepackage{hyperref}
\usepackage{lastpage}
\usepackage{amsmath}
\usepackage{layout}

\lhead{}
\chead{}
%%%%%%%%%%%%%%%%%%%%%%%%%%%%%%%%%%%%%%%%%%%%%%%%%%%%%%%%%%%%%%

% Modify header here %%%%%%%%%%%%%%%%%%%%%%%%%%%%%%%%%%%%%%%%%
\rhead{\footnotesize Microeconom\'ia Aplicada, Oto\~no 2025}

%%%%%%%%%%%%%%%%%%%%%%%%%%%%%%%%%%%%%%%%%%%%%%%%%%%%%%%%%%%%%%
% Don't touch this %%%%%%%%%%%%%%%%%%%%%%%%%%%%%%%%%%%%%%%%%%%
\lfoot{}
\cfoot{\small \thepage/\pageref*{LastPage}}
\rfoot{}

% \usepackage{array, xcolor}
% \usepackage{color,hyperref}
% \definecolor{clemsonorange}{HTML}{EA6A20}
% \hypersetup{colorlinks,breaklinks,linkcolor=clemsonorange,urlcolor=clemsonorange,anchorcolor=clemsonorange,citecolor=black}

%===================Startup========================
\begin{document}

\maketitle

\blankline

\setlength{\parindent}{0pt}.

1.- (Negociaciones de salario) Este es un juego de negociaciones entre un sindicato laboral y una firma. En una primera etapa, el sindicator decide que salario debe recibir $w$, el cual afecta a todos los trabajadores. En la segunda etapa, la firma decide cuantos trabajadores contratar $h \in [0,1]$. La funcion de utilidad del sindicato es
\begin{align*}
	u_{L}(w,h) = wh
\end{align*} por lo que se incrementa con ambos argumentos. De otro modo, la de la firma es
\begin{align*}
	\pi(w,h) = \left( h - \frac{h^{2}}{2} \right) - wh
\end{align*}

\begin{enumerate}
	\item Obtenga el/los equilibrio de Nash perfecto en subjuegos.
	\item ¿Cu\'al es la historia terminal asociada al equilibrio obtenido en la pregunta anterior?
	\item ¿Cu\'ales son los pagos asociados al equilibrio obtenido en la primera pregunta?
\end{enumerate}

2.- (Stackelberg con 3 firmas) Una industria tiene 3 firmas ($i=1,2,3$). Cada firma $i$ tiene la misma estructura de costos, la cual esta dada por $C(q_{i}) = 5 + 2q_{i}$, donde $q_{i}$ es la cantidad producida por firma $i$. Adicionalmente, la funcion inversa de demanda esta dada por $P(Q) = 18 - Q$. Asuma que primero la firma 1 produce, luego la 2 y finalmente la 3.

\begin{enumerate}
	\item Obtenga el/los equilibrio de Nash perfecto en subjuegos.
	\item Compare el excedente del consumidor de la pregunta anterior con el resultante de un juego de Cournot con 3 firmas (con la misma estructura de costos).
\end{enumerate}

3.- (Bloqueo de entrada) Existen 2 firmas. La incumbente (firma 1) tiene que decidir si invierte o no en una tecnologia. La inversion es observada por el potencial entrante (firma 2) que decide si entra o no. Si la firma 2 no entra, el juego termina y la firma 1 se comporta como un monopolista. Si la firma 2 entra, las firmas juegan simultaneamente decidiendo precios bajos (L), medios (M) o altos (H) segun la figura anterior.
\begin{figure}[ht!]
\centering
\includegraphics[width=0.8\textwidth]{../../../src/original_data/ayudantia2.jpg}
%\caption{A simple caption \label{overflow}}
\end{figure}


\begin{enumerate}
	\item Obtenga el/los equilibrio de Nash perfecto en subjuegos.
\end{enumerate}

\end{document}  